\section{Geometric input and normalization}\label{sec:geometry}
We use established regularity and structure theorems for surface mean curvature flow in their stated scope. We record precisely the consequences required below.

The endpoint and test clause in the following lemma is the integrated
one used in \cite[Section~2.1]{BK2024}.  We state separately the local
integrability assumptions needed for compact tests.

\begin{lemma}[Equality of all-endpoint Brakke representatives]
\label{geo:representative}
Let $I\subset(-\infty,0)$ be an open interval, let $S$ be a smooth
properly embedded self-shrinker without boundary, and put
$\nu_s=\mathcal H^2\!\llcorner(\sqrt{-s}\,S)$.  Let $(\mu_s)_{s\in I}$
be a weakly measurable family of Radon measures.  Suppose that for
almost every $s$, $\mu_s$ is the weight of a rectifiable $2$-varifold
$V_s$ with measurable generalized mean curvature $\mathbf H_\mu$ and
first variation $\delta V_s=-\mathbf H_\mu\mu_s$ on compactly
supported vector fields.  Suppose also that for every compact
$K\subset\mathbb R^3$ and compact interval $J\Subset I$,
\begin{equation}\label{geo:representative-integrability}
 \int_J\!\int_K(1+|\mathbf H_\mu|^2)\,d\mu_s\,ds<\infty.
\end{equation}
Assume that for every $a<b$ in $I$ and every nonnegative
$\varphi\in C_c^1(\mathbb R^3\times[a,b])$,
\begin{multline}\label{geo:all-endpoint-brakke}
 \mu_b(\varphi(\cdot,b))-\mu_a(\varphi(\cdot,a))
 \\
 \le\int_a^b\!\int\bigl(\partial_s\varphi
       +\nabla\varphi\cdot\mathbf H_\mu
       -\varphi|\mathbf H_\mu|^2\bigr)\,d\mu_s\,ds.
\end{multline}
The test is defined on the closed slab and may be nonzero at its
temporal endpoints, and $\partial_s$ is taken at fixed ambient
position.  Both $a$ and $b$ are arbitrary, including exceptional
times.  If $\mu_s=\nu_s$ as Radon measures for almost every $s\in I$,
then $\mu_s=\nu_s$ for every $s\in I$.
\end{lemma}

\begin{proof}
The local mass and curvature bounds in
\eqref{geo:representative-integrability} make every compact-test
integrand in \eqref{geo:all-endpoint-brakke} locally integrable;
the mixed term follows by Cauchy--Schwarz.  The smooth reference has
the analogous local bounds and the compact-test transport identity.
Indeed, if $J\Subset(-\infty,0)$ is a compact interval and a compact
spatial set lies in $B_R$, the preimage of its spacetime track under
$(p,s)\mapsto(\sqrt{-s}\,p,s)$ lies in
\[
 \bigl(S\cap\overline B_{R/\min_{s\in J}\sqrt{-s}}\bigr)\times J.
\]
This set is compact by properness of $S$.  Smoothness on that compact
set bounds the local area and curvature integrals and justifies the
usual transport formula for compact tests on $J$.  No finite global
unweighted area or global integrability of $\mathbf H$ is required.

At almost every time where $\mu_s=\nu_s$, their associated rectifiable
varifolds also agree: a rectifiable weight determines its approximate
tangent plane almost everywhere.  Their first variations, and hence
their generalized mean curvature densities, agree at those times.
Write $\mathbf H_\nu$ for the mean curvature vector of the smooth
surface $\sqrt{-s}\,S$.  Fix a nonnegative
$\phi\in C_c^2(\mathbb R^3)$ and write
\[
 \beta_\phi(s)=\int\bigl(\nabla\phi\cdot\mathbf H_\nu
                    -\phi|\mathbf H_\nu|^2\bigr)\,d\nu_s.
\]
The preceding local transport identity gives, for every $a<b$ in $I$,
\[
 \nu_b(\phi)-\nu_a(\phi)=\int_a^b\beta_\phi(s)\,ds.
\]
Apply \eqref{geo:all-endpoint-brakke} to the time-independent test
$\varphi(X,s)=\phi(X)$.  Its integrand equals $\beta_\phi$ almost
everywhere, so
\[
 \mu_b(\phi)-\mu_a(\phi)
 \le\nu_b(\phi)-\nu_a(\phi)\qquad(a<b).
\]
Thus $h_\phi(s)=\mu_s(\phi)-\nu_s(\phi)$ is a finite nonincreasing
function and vanishes almost everywhere.  For any $s\in I$, choose
times $a<s<b$ in the full-measure set where $h_\phi=0$.  Then
\[
 0=h_\phi(a)\ge h_\phi(s)\ge h_\phi(b)=0.
\]
Nonnegative smooth compact tests determine Radon measures, proving
the assertion.  The two-sided choice of good times is why the lemma
asserts equality at interior times; every $s<0$ is interior when
$I=(-\infty,0)$.
\end{proof}


\begin{proposition}[Tangent section and ends]\label{geo:tangent}
A backward tangent flow of the flow in \Cref{thm:main} has multiplicity one and is the homothetic motion $\sqrt{-s}\,S$, $s<0$, of a smooth properly embedded self-shrinking surface. Its time slices have uniform quadratic local area bounds. The surface $S$ has bounded geometry and finite genus. Outside a compact set it is the disjoint union of finitely many ends, each either smoothly asymptotic to a regular cone under dilations or a cylindrical tail whose translates centered at points receding outward along a ray from the origin converge smoothly on compact sets to the full round cylinder of radius $\sqrt2$. The corresponding links and cylindrical directions are separated on the unit sphere.
\end{proposition}

\begin{proof}
The smooth flow with compact slices lies in the almost-regular-flow
framework of Bamler--Kleiner by their Theorem~1.1.  It is bounded in
their terminology, and their Theorems~1.2 and~1.7(b) give the applicable
multiplicity-one and compactness statements, including tangents at
the finite endpoint of the smooth time interval \cite{BK2024}.
We use the v3 formulations.  Fix the scales extracting a tangent in the
convention of \Cref{sec:introduction}.  On every fixed interval
$(a,b)\Subset(-\infty,0)$ these rescalings have uniform entropy and
genus bounds.  Apply Theorem~1.7(b) first to extract, on a subsequence,
an associated almost-regular limit.  The Gaussian density drop between the endpoints tends
to zero, since both endpoint densities tend by monotonicity to the
same density at $(x_0,T)$.  Lemma~2.13(b) of \cite{BK2024}, including
its genus conclusion in equation~(2.16), therefore characterizes the
weights of that limit as those of a smooth properly embedded homothetic
shrinker of finite genus on the interval; Theorem~1.2 supplies multiplicity one.
By Theorem~1.7(b), this same subsequence converges locally
smoothly at every regular time of its associated almost-regular limit.
Those times have full measure by Definition~2.6(1), and Lemma~2.7
identifies the associated slice weights with the areas of their regular
surfaces.  Hence the slice weights converge locally weakly for almost
every time.  The uniform local area bounds and dominated convergence
give the local spacetime weight of this smooth homothetic flow.  The original extracted scales
already have the spacetime weight of the specified tangent.  Uniqueness
of the local weak Radon limit identifies those spacetime measures.
Testing by a countable determining family of smooth compact spatial
functions times arbitrary compact time functions then identifies the
slice weights for almost every time.  Lemma~\ref{geo:representative}
gives equality at every interior time of the interval.  The homothetic
descriptions therefore agree on overlapping intervals and assemble to
the asserted motion on all $s<0$.
Applying the same compactness and density-drop argument to any scale
sequence on an exhaustion of $(-\infty,0)$, taking a diagonal
subsequence, and using the resulting smooth homothetic flow as its
representative, also shows that the tangent class in the stated
convention is nonempty: smoothness supplies its local curvature and
all-endpoint transport conditions, and the preceding convergence gives
its spacetime weight.
Their Theorem~1.10 gives the finite conical/cylindrical end decomposition.
Huisken's monotonicity formula supplies quadratic area ratios from the
finite entropy of the compact initial surface \cite{Huisken1990};
the density monotonicity statement is also recorded in
\cite[Section~2.9]{White2005}.  Bounded geometry on fixed spatial
scales follows from the smooth plane or cylinder limits at points
escaping to infinity in the proof of the end theorem
\cite[Claim~7.6]{BK2024}, together with the smooth compact core.
The higher end estimates in logarithmic coordinates used below will
be derived from the stationary equation.
\end{proof}

Lu Wang's end theorem gives the conical/cylindrical alternative for noncompact properly embedded self-shrinkers in $\R^3$ of finite topology \cite[Theorem~1.1]{Wang2016}. Here the Bamler--Kleiner theorem in \Cref{geo:tangent} supplies the end decomposition directly for the tangent class; finite genus alone would not imply Wang's finite-topology hypothesis. We use neither a rigidity theorem for an arbitrary shrinker with a cylindrical end nor the tangent-flow uniqueness that is to be proved.

Translate $(x_0,T)$ to $(0,0)$ and write the original flow on $t<0$. Its rescaled form is
\begin{equation}\label{geo:rescaled}
 M(\tau)=e^{\tau/2}M_{-e^{-\tau}}.
\end{equation}
We omit the inessential constant prefactor in the Gaussian area and put
\begin{equation}\label{geo:action}
 F(M)=\int_M e^{-|X|^2/4}\dd\mu,
 \qquad \Phi=\mathbf H+\tfrac12X^\perp,
 \qquad d(\tau)^2=\int_{M(\tau)}|\Phi|^2e^{-|X|^2/4}\dd\mu.
\end{equation}
Then $-\partial_\tau F(M(\tau))=d(\tau)^2$. Fix one tangent section $S$ from \Cref{geo:tangent}. Along a corresponding subsequence, convergence is smooth on compact subsets in the interior of negative time. Indeed, the approximating smooth integral Brakke flows have bounded area ratios and converge locally in measure to a smooth multiplicity-one flow, so the local Brakke-regularity consequence recorded in \cite[Lemma~2.1]{Schulze2014} applies after shrinking each compact spacetime cylinder; see also \cite{White2005}. Quadratic area bounds make Gaussian tails uniformly small, so
\begin{equation}\label{geo:energy}
 \mathcal E(\tau)=F(M(\tau))-F(S)\searrow0.
\end{equation}
The letter $E$ below is reserved for the small coefficient of radial diffusion, and is distinct from the energy $\mathcal E$.

For complex stationary graphs, $X^2$ means the complex bilinear product $X\cdot X$, without conjugation. Metric inverses, curvature, and area elements are complexified in the same way. Branches are continued from the positive real area element. The independent spatial parameters remain real; no holomorphic continuation of an arbitrary smooth spatial graph is assumed. The compact observations used in the normalized stationary equation are kept in fixed real graph charts.

All constants associated with one fixed end list, a fixed angular regularity order, or a fixed contour-size parameter may depend on those choices. Constants asserted uniform in an opening scale do not. A small loss in a spatial radius may be divided among finitely many cutoffs; this convention will always be used after the relevant exponents have been fixed.
